% Options for packages loaded elsewhere
\PassOptionsToPackage{unicode}{hyperref}
\PassOptionsToPackage{hyphens}{url}
\PassOptionsToPackage{dvipsnames,svgnames,x11names}{xcolor}
%
\documentclass[
]{article}
\usepackage{amsmath,amssymb}
\usepackage{iftex}
\ifPDFTeX
  \usepackage[T1]{fontenc}
  \usepackage[utf8]{inputenc}
  \usepackage{textcomp} % provide euro and other symbols
\else % if luatex or xetex
  \usepackage{unicode-math} % this also loads fontspec
  \defaultfontfeatures{Scale=MatchLowercase}
  \defaultfontfeatures[\rmfamily]{Ligatures=TeX,Scale=1}
\fi
\usepackage{lmodern}
\ifPDFTeX\else
  % xetex/luatex font selection
  \setmainfont[]{Inter}
\fi
% Use upquote if available, for straight quotes in verbatim environments
\IfFileExists{upquote.sty}{\usepackage{upquote}}{}
\IfFileExists{microtype.sty}{% use microtype if available
  \usepackage[]{microtype}
  \UseMicrotypeSet[protrusion]{basicmath} % disable protrusion for tt fonts
}{}
\makeatletter
\@ifundefined{KOMAClassName}{% if non-KOMA class
  \IfFileExists{parskip.sty}{%
    \usepackage{parskip}
  }{% else
    \setlength{\parindent}{0pt}
    \setlength{\parskip}{6pt plus 2pt minus 1pt}}
}{% if KOMA class
  \KOMAoptions{parskip=half}}
\makeatother
\usepackage{xcolor}
\usepackage[margin=18mm]{geometry}
\usepackage{longtable,booktabs,array}
\usepackage{calc} % for calculating minipage widths
% Correct order of tables after \paragraph or \subparagraph
\usepackage{etoolbox}
\makeatletter
\patchcmd\longtable{\par}{\if@noskipsec\mbox{}\fi\par}{}{}
\makeatother
% Allow footnotes in longtable head/foot
\IfFileExists{footnotehyper.sty}{\usepackage{footnotehyper}}{\usepackage{footnote}}
\makesavenoteenv{longtable}
\setlength{\emergencystretch}{3em} % prevent overfull lines
\providecommand{\tightlist}{%
  \setlength{\itemsep}{0pt}\setlength{\parskip}{0pt}}
\setcounter{secnumdepth}{-\maxdimen} % remove section numbering
\ifLuaTeX
  \usepackage{selnolig}  % disable illegal ligatures
\fi
\IfFileExists{bookmark.sty}{\usepackage{bookmark}}{\usepackage{hyperref}}
\IfFileExists{xurl.sty}{\usepackage{xurl}}{} % add URL line breaks if available
\urlstyle{same}
\hypersetup{
  colorlinks=true,
  linkcolor={blue},
  filecolor={Maroon},
  citecolor={Blue},
  urlcolor={Blue},
  pdfcreator={LaTeX via pandoc}}

\author{}
\date{}

\begin{document}

\hypertarget{documento-de-validaciuxf3n-con-contraparte--plataforma-cicuc}{%
\section{Documento de validación con contraparte --- Plataforma
CICUC}\label{documento-de-validaciuxf3n-con-contraparte--plataforma-cicuc}}

\textbf{Versión:} 1.0\\
\textbf{Fecha:} 31 de julio de 2026\\
\textbf{Estado:} Propuesta para revisión clínica, operativa y
administrativa\\
\textbf{Proyecto:} Plataforma de gestión de estudios clínicos
oncológicos CICUC

\begin{quote}
Este documento busca confirmar procesos, alcance y responsabilidades
antes de continuar el desarrollo. No autoriza el tratamiento de datos
clínicos reales, no reemplaza los protocolos oficiales y no permite que
la plataforma determine elegibilidad clínica.
\end{quote}

\hypertarget{1-objetivo-de-la-validaciuxf3n}{%
\subsection{1. Objetivo de la
validación}\label{1-objetivo-de-la-validaciuxf3n}}

Confirmar con CICUC que el equipo de desarrollo comprendió
correctamente:

\begin{itemize}
\tightlist
\item
  el problema operacional;
\item
  los usuarios y responsables;
\item
  el ciclo de vida de los estudios;
\item
  la gestión de slots, listas de espera y pretesting;
\item
  el alcance del primer producto funcional;
\item
  las decisiones institucionales pendientes.
\end{itemize}

Los acuerdos obtenidos se incorporarán a las historias de usuario,
criterios de aceptación, diseño y pruebas.

\hypertarget{2-problema-comprendido}{%
\subsection{2. Problema comprendido}\label{2-problema-comprendido}}

Actualmente, la información sobre estudios puede encontrarse distribuida
entre PDF, Word, correos y notas. Esto dificulta saber con certeza:

\begin{itemize}
\tightlist
\item
  qué estudios están disponibles;
\item
  cuál es su estado;
\item
  qué protocolo está vigente;
\item
  quién debe actualizar la información;
\item
  qué slots, pretests y tareas están pendientes;
\item
  qué compromisos de reclutamiento requieren seguimiento.
\end{itemize}

La necesidad prioritaria es una fuente de verdad operacional,
actualizada, gobernada y trazable.

\hypertarget{3-resultado-propuesto-para-el-mvp}{%
\subsection{3. Resultado propuesto para el
MVP}\label{3-resultado-propuesto-para-el-mvp}}

Un usuario autorizado podrá:

\begin{enumerate}
\def\labelenumi{\arabic{enumi}.}
\tightlist
\item
  consultar estudios por patología, línea, biomarcador, centro y estado;
\item
  conocer la versión vigente del protocolo y quién la confirmó;
\item
  distinguir preactivación, reclutamiento, seguimiento, hold, cierre y
  archivo;
\item
  consultar cohortes, ramas y criterios cargados manualmente;
\item
  gestionar slots, lista de espera, pretesting y tareas;
\item
  consultar reportes operativos con fecha de corte;
\item
  reconstruir cambios relevantes mediante auditoría.
\item
  disponer de una arquitectura modular y APIs documentadas, preparada
  para integraciones graduales con sistemas institucionales autorizados.
\end{enumerate}

\hypertarget{4-alcance-que-no-se-propone-para-el-mvp}{%
\subsection{4. Alcance que no se propone para el
MVP}\label{4-alcance-que-no-se-propone-para-el-mvp}}

\begin{itemize}
\tightlist
\item
  determinación automática de elegibilidad;
\item
  extracción automática de protocolos;
\item
  análisis de audio o escucha de comités;
\item
  integración con ficha clínica;
\item
  mensajería externa;
\item
  recomendación o priorización clínica automatizada.
\end{itemize}

Estas capacidades podrán evaluarse posteriormente con validación humana,
evidencia trazable, controles de seguridad y aprobación institucional.

La preparación para interoperabilidad no significa que el MVP se
conectará automáticamente con fichas clínicas u otros sistemas. Cada
integración futura requerirá alcance, contrato, permisos, seguridad y
pruebas específicas.

\hypertarget{5-plan-de-ejecuciuxf3n-propuesto}{%
\subsection{5. Plan de ejecución
propuesto}\label{5-plan-de-ejecuciuxf3n-propuesto}}

\begin{itemize}
\tightlist
\item
  \textbf{Mes 1:} levantamiento de información, requerimientos y flujo
  de procesos/arquitectura.
\item
  \textbf{Mes 2:} desarrollo de MVP funcional básico y validación con
  usuarios.
\item
  \textbf{Mes 3:} puesta en marcha y marcha blanca.
\end{itemize}

\hypertarget{6-ciclo-de-vida-propuesto}{%
\subsection{6. Ciclo de vida
propuesto}\label{6-ciclo-de-vida-propuesto}}

\begin{longtable}[]{@{}ll@{}}
\toprule\noalign{}
Estado & Definición propuesta \\
\midrule\noalign{}
\endhead
\bottomrule\noalign{}
\endlastfoot
Preactivación & Estudio próximo a abrir, potencialmente dentro de unas
cuatro semanas \\
Activo reclutando & Estudio abierto y buscando participantes \\
Activo en seguimiento & No recluta, pero mantiene participantes en
control \\
Hold & Reclutamiento temporalmente detenido \\
Cerrado & Sin reclutamiento ni seguimiento operativo activo \\
Archivado & Conservado únicamente para trazabilidad e historia \\
\end{longtable}

El estado del estudio debe permanecer separado de la disponibilidad de
cupos o slots.

\hypertarget{7-informaciuxf3n-muxednima-del-inventario}{%
\subsection{7. Información mínima del
inventario}\label{7-informaciuxf3n-muxednima-del-inventario}}

\begin{itemize}
\tightlist
\item
  código, nombre corto y título oficial;
\item
  patrocinador, fase y diseño;
\item
  patología, subtipo, escenario y línea;
\item
  cohortes, ramas y biomarcadores;
\item
  centro, investigador y coordinadora responsable;
\item
  estado operativo;
\item
  disponibilidad de cupo o slot;
\item
  versión vigente del protocolo;
\item
  fecha, fuente y responsable de la última verificación;
\item
  próxima fecha de revisión.
\end{itemize}

\hypertarget{8-flujo-propuesto-de-slots}{%
\subsection{8. Flujo propuesto de
slots}\label{8-flujo-propuesto-de-slots}}

\begin{enumerate}
\def\labelenumi{\arabic{enumi}.}
\tightlist
\item
  Identificar anticipadamente una oportunidad potencial.
\item
  Solicitar un slot o registrar ingreso a lista de espera.
\item
  Registrar una prioridad operativa justificada.
\item
  Registrar confirmación del patrocinador.
\item
  Mostrar la ventana disponible y sus fechas.
\item
  Revisar criterios y datos dinámicos de manera humana.
\item
  Registrar utilización, liberación o vencimiento.
\end{enumerate}

La lista de espera no constituye un ranking clínico. Cada movimiento
debe registrar responsable, fecha, fuente y motivo.

\hypertarget{9-flujo-propuesto-de-pretesting}{%
\subsection{9. Flujo propuesto de
pretesting}\label{9-flujo-propuesto-de-pretesting}}

La plataforma permitirá seguir:

\begin{itemize}
\tightlist
\item
  prueba requerida;
\item
  estudio o cohorte relacionada;
\item
  fecha de indicación y toma;
\item
  estado del resultado;
\item
  fuente;
\item
  vigencia o vencimiento;
\item
  responsable y próxima acción;
\item
  condición que gatillaría una revisión.
\end{itemize}

Los recordatorios serán administrativos y no modificarán automáticamente
una preselección ni reservarán cupos.

\hypertarget{10-usuarios-propuestos}{%
\subsection{10. Usuarios propuestos}\label{10-usuarios-propuestos}}

\begin{itemize}
\tightlist
\item
  investigador principal;
\item
  médico investigador u oncólogo derivador;
\item
  enfermera coordinadora;
\item
  jefatura de enfermería;
\item
  coordinación administrativa;
\item
  equipos de factibilidad y preactivación;
\item
  administración técnica;
\item
  auditor autorizado.
\end{itemize}

Los permisos se definirán según centro, estudio, acción y relación con
el recurso. Un cargo o participación en una reunión no concederá
automáticamente acceso clínico.

\hypertarget{11-datos-y-resguardos}{%
\subsection{11. Datos y resguardos}\label{11-datos-y-resguardos}}

Antes de utilizar información real, CICUC deberá aprobar:

\begin{itemize}
\tightlist
\item
  datos mínimos permitidos;
\item
  fundamento o consentimiento aplicable;
\item
  permisos de visualización y modificación;
\item
  separación de información identificatoria, clínica y operativa;
\item
  retención y eliminación;
\item
  custodia de protocolos y documentos;
\item
  reglas de descarga y exportación;
\item
  auditoría y respuesta ante incidentes.
\end{itemize}

Hasta esa aprobación, el desarrollo y las demostraciones utilizarán
únicamente datos ficticios.

\hypertarget{12-distribuciuxf3n-propuesta-de-responsabilidades-sobre-datos}{%
\subsection{12. Distribución propuesta de responsabilidades sobre
datos}\label{12-distribuciuxf3n-propuesta-de-responsabilidades-sobre-datos}}

Esta distribución debe incorporarse a un contrato y a un anexo de
tratamiento de datos revisados por asesoría jurídica.

\hypertarget{cicuc-como-responsable-de-los-datos}{%
\subsubsection{CICUC como responsable de los
datos}\label{cicuc-como-responsable-de-los-datos}}

CICUC, como institución que determina los fines y medios del
tratamiento, debería asumir la responsabilidad por:

\begin{itemize}
\tightlist
\item
  definir la finalidad y base de licitud de cada tratamiento;
\item
  determinar qué datos pueden incorporarse y su plazo de conservación;
\item
  asegurar la legitimidad, exactitud y procedencia de los datos
  entregados;
\item
  obtener consentimientos o autorizaciones cuando correspondan;
\item
  gestionar derechos de los titulares y comunicaciones institucionales;
\item
  autorizar usuarios, perfiles, integraciones, exportaciones y
  subencargados;
\item
  custodiar las fichas clínicas y documentos fuente;
\item
  adoptar las decisiones clínicas y operacionales basadas en la
  información.
\end{itemize}

\hypertarget{adyac-como-encargado-o-proveedor-tecnoluxf3gico}{%
\subsubsection{ADYAC como encargado o proveedor
tecnológico}\label{adyac-como-encargado-o-proveedor-tecnoluxf3gico}}

ADYAC trataría datos únicamente por cuenta de CICUC y conforme a
instrucciones documentadas. Sus obligaciones deberían limitarse a:

\begin{itemize}
\tightlist
\item
  construir y operar los componentes contratados;
\item
  aplicar controles técnicos y organizativos proporcionales al riesgo;
\item
  mantener confidencialidad y mínimo privilegio;
\item
  no utilizar datos para finalidades propias;
\item
  informar incidentes o instrucciones que parezcan contrarias al
  contrato;
\item
  colaborar en auditorías y solicitudes dentro del alcance acordado;
\item
  devolver o eliminar datos al finalizar, según instrucciones y
  obligaciones legales aplicables.
\end{itemize}

ADYAC no debería asumir responsabilidad por la legalidad de los datos
seleccionados por CICUC, la falta de consentimiento o autorización
institucional, la inexactitud de antecedentes ingresados por usuarios,
el uso fuera del alcance contratado ni decisiones clínicas u
operacionales tomadas por CICUC. Esto no excluye la responsabilidad
propia de ADYAC por incumplimiento contractual, dolo, culpa grave o
vulneración de sus deberes legales directos.

\hypertarget{13-mecanismos-contractuales-y-de-protecciuxf3n-recomendados}{%
\subsection{13. Mecanismos contractuales y de protección
recomendados}\label{13-mecanismos-contractuales-y-de-protecciuxf3n-recomendados}}

Antes de producción se recomienda formalizar:

\begin{enumerate}
\def\labelenumi{\arabic{enumi}.}
\tightlist
\item
  \textbf{Contrato principal y alcance:} entregables, exclusiones,
  dependencias, criterios de aceptación y control de cambios.
\item
  \textbf{Anexo de tratamiento de datos:} roles, instrucciones,
  categorías de datos, finalidades, retención, devolución, eliminación y
  asistencia.
\item
  \textbf{Anexo de seguridad:} cifrado, accesos, registros, respaldo,
  vulnerabilidades, continuidad, ambientes y evidencia de pruebas.
\item
  \textbf{Matriz de responsabilidades:} quién autoriza usuarios, corrige
  datos, responde solicitudes, comunica incidentes y decide
  restauraciones.
\item
  \textbf{Subencargados:} autorización previa, listado de proveedores,
  ubicación de datos y obligaciones equivalentes.
\item
  \textbf{Incidentes:} canal, plazos de aviso, preservación de
  evidencia, investigación y coordinación de comunicaciones.
\item
  \textbf{Indemnidad recíproca:} CICUC responde por reclamos derivados
  de datos, instrucciones o decisiones bajo su control; ADYAC por
  incumplimientos atribuibles a sus propias obligaciones.
\item
  \textbf{Límite de responsabilidad:} monto máximo negociado, exclusión
  de daños indirectos y excepciones para dolo, culpa grave y materias
  que la ley no permita limitar.
\item
  \textbf{Propiedad intelectual:} titularidad del desarrollo específico,
  componentes preexistentes, licencias, software de terceros y
  continuidad de uso.
\item
  \textbf{Aceptación y garantía:} pruebas, plazo de observaciones,
  defectos cubiertos y diferencia entre corrección y nueva
  funcionalidad.
\item
  \textbf{Seguros:} evaluar responsabilidad civil profesional,
  ciberseguridad y cobertura de incidentes de acuerdo con el riesgo
  real.
\item
  \textbf{Solución de controversias:} ley aplicable, jurisdicción,
  notificación y mecanismo previo de escalamiento y negociación.
\end{enumerate}

Ninguna cláusula puede impedir que un tercero presente una reclamación.
El objetivo es asignar responsabilidades de forma clara, reducir la
probabilidad de incidentes y definir quién soporta contractualmente sus
consecuencias.

\hypertarget{14-marco-legal-de-referencia}{%
\subsection{14. Marco legal de
referencia}\label{14-marco-legal-de-referencia}}

\begin{itemize}
\tightlist
\item
  Ley Nº 19.628, sobre protección de la vida privada, mientras continúe
  vigente.
\item
  Ley Nº 21.719, que moderniza la protección de datos y entra en
  vigencia el 1 de diciembre de 2026. Distingue al responsable que
  decide fines y medios y al encargado que trata datos por su cuenta,
  exigiendo medidas adecuadas al riesgo.
\item
  Ley Nº 20.584: la información de ficha clínica y estudios asociados es
  dato sensible; su custodia, acceso, confidencialidad e
  interoperabilidad están sujetas a reglas sectoriales.
\item
  Ley Nº 21.663, Marco de Ciberseguridad: contempla medidas permanentes
  para prevenir, reportar y resolver incidentes en las instituciones
  alcanzadas.
\item
  Ley Nº 17.336: protege programas computacionales y documentación; el
  software desarrollado por encargo se presume cedido al cliente salvo
  pacto escrito en contrario.
\item
  Código Civil, artículo 1558: permite modificar contractualmente
  determinadas reglas de responsabilidad, sin convertir una exención
  total en garantía de inmunidad.
\end{itemize}

Esta sección es una propuesta de gestión contractual, no un informe
jurídico. El contrato definitivo debe ser revisado por abogados de ambas
partes.

\hypertarget{15-reportes-propuestos}{%
\subsection{15. Reportes propuestos}\label{15-reportes-propuestos}}

\begin{itemize}
\tightlist
\item
  estudios por estado, patología y centro;
\item
  reclutamiento comprometido versus real;
\item
  slots solicitados, asignados, utilizados y vencidos;
\item
  lista de espera;
\item
  pretests y tareas pendientes;
\item
  antigüedad de la información;
\item
  actividad y trazabilidad operacional.
\end{itemize}

Cada indicador deberá contar con definición, fuente, periodo, fecha de
corte, unidad de análisis y permisos aprobados.

\hypertarget{16-decisiones-solicitadas-a-cicuc}{%
\subsection{16. Decisiones solicitadas a
CICUC}\label{16-decisiones-solicitadas-a-cicuc}}

\begin{longtable}[]{@{}lll@{}}
\toprule\noalign{}
Tema & Propuesta & Validación \\
\midrule\noalign{}
\endhead
\bottomrule\noalign{}
\endlastfoot
Estados del estudio & Preactivación, reclutamiento, seguimiento, hold,
cerrado y archivado & ☐ Aprobar ☐ Ajustar \\
Responsable del inventario & Coordinadora por estudio con supervisión de
jefatura & ☐ Aprobar ☐ Ajustar \\
Vigencia & Cada registro muestra fecha, fuente, responsable y próxima
revisión & ☐ Aprobar ☐ Ajustar \\
Slots & Solicitud, espera, asignación, uso, liberación y vencimiento & ☐
Aprobar ☐ Ajustar \\
Lista de espera & Prioridad operacional trazable, no ranking clínico & ☐
Aprobar ☐ Ajustar \\
Pretesting & Tareas y recordatorios administrativos con fechas y
responsables & ☐ Aprobar ☐ Ajustar \\
Pacientes & Solo datos ficticios hasta aprobar privacidad y permisos & ☐
Aprobar ☐ Ajustar \\
Elegibilidad & Siempre corresponde al investigador y equipo clínico & ☐
Aprobar ☐ Ajustar \\
IA y voz & Fuera del MVP administrativo & ☐ Aprobar ☐ Ajustar \\
Reportes & Implementar después de aprobar definiciones de KPI & ☐
Aprobar ☐ Ajustar \\
Responsabilidad de datos & CICUC como responsable; ADYAC como encargado
bajo instrucciones & ☐ Aprobar ☐ Ajustar \\
Contrato de datos & Suscribir anexo de tratamiento y seguridad antes de
producción & ☐ Aprobar ☐ Ajustar \\
\end{longtable}

\hypertarget{17-preguntas-para-la-reuniuxf3n}{%
\subsection{17. Preguntas para la
reunión}\label{17-preguntas-para-la-reuniuxf3n}}

\begin{enumerate}
\def\labelenumi{\arabic{enumi}.}
\tightlist
\item
  ¿Quién mantiene cada estudio y quién lo reemplaza?
\item
  ¿Cada cuánto debe confirmarse la información?
\item
  ¿Cuándo un registro debe mostrarse como desactualizado?
\item
  ¿Qué estados y transiciones se usan actualmente?
\item
  ¿Qué diferencia operacional existe entre cupo y slot?
\item
  ¿Qué datos mínimos necesita una lista de espera?
\item
  ¿Qué pretests se anticipan y cuánto dura su vigencia?
\item
  ¿Dónde se custodiarán protocolos confidenciales?
\item
  ¿Qué puede consultar un oncólogo externo?
\item
  ¿Cómo se definen los compromisos y KPI de reclutamiento?
\item
  ¿Quién actuará formalmente como responsable de datos dentro de CICUC?
\item
  ¿Qué proveedores o subencargados institucionales están autorizados?
\item
  ¿Qué procedimiento institucional se aplicará ante incidentes?
\end{enumerate}

\hypertarget{18-registro-de-validaciuxf3n}{%
\subsection{18. Registro de
validación}\label{18-registro-de-validaciuxf3n}}

\textbf{Fecha de reunión:}
\_\_\_\_\_\_\_\_\_\_\_\_\_\_\_\_\_\_\_\_\_\_\_\_\_\_\_\_\_\_\_\_\_\_\_\_\_\_\_\_\_\_\_\_\_\_\_\_\_

\textbf{Participantes:}
\_\_\_\_\_\_\_\_\_\_\_\_\_\_\_\_\_\_\_\_\_\_\_\_\_\_\_\_\_\_\_\_\_\_\_\_\_\_\_\_\_\_\_\_\_\_\_\_\_\_\_\_

\textbf{Acuerdos principales:}

\begin{center}\rule{0.5\linewidth}{0.5pt}\end{center}

\begin{center}\rule{0.5\linewidth}{0.5pt}\end{center}

\textbf{Observaciones o cambios requeridos:}

\begin{center}\rule{0.5\linewidth}{0.5pt}\end{center}

\begin{center}\rule{0.5\linewidth}{0.5pt}\end{center}

\begin{longtable}[]{@{}ll@{}}
\toprule\noalign{}
Por CICUC & Por el equipo del proyecto \\
\midrule\noalign{}
\endhead
\bottomrule\noalign{}
\endlastfoot
Nombre: \_\_\_\_\_\_\_\_\_\_\_\_\_\_\_\_\_\_\_\_\_\_\_\_\_\_ & Nombre:
\_\_\_\_\_\_\_\_\_\_\_\_\_\_\_\_\_\_\_\_\_\_\_\_\_\_ \\
Cargo: \_\_\_\_\_\_\_\_\_\_\_\_\_\_\_\_\_\_\_\_\_\_\_\_\_\_\_ & Cargo:
\_\_\_\_\_\_\_\_\_\_\_\_\_\_\_\_\_\_\_\_\_\_\_\_\_\_\_ \\
Firma: \_\_\_\_\_\_\_\_\_\_\_\_\_\_\_\_\_\_\_\_\_\_\_\_\_\_\_ & Firma:
\_\_\_\_\_\_\_\_\_\_\_\_\_\_\_\_\_\_\_\_\_\_\_\_\_\_\_ \\
Fecha: \_\_\_\_\_\_\_\_\_\_\_\_\_\_\_\_\_\_\_\_\_\_\_\_\_\_\_ & Fecha:
\_\_\_\_\_\_\_\_\_\_\_\_\_\_\_\_\_\_\_\_\_\_\_\_\_\_\_ \\
\end{longtable}

\end{document}
